\subsubsection{Sensibilité de la topologie à la largeur de bande}\label{sec:sensibilite-arbres}
Une bande de demi-largeur \(b\) conserve les cellules \(|i-j|\le b\). Si \(|m-n|>b\), le point d’arrivée est inaccessible. Sinon, le score contraint vérifie \(d_b\ge d\). Edlib retrouve les 2\,016 distances publiées pour les 64 génomes du SARS-CoV-2. Si son chemin tient dans la bande, \(d_b=d\). Pour les autres paires, la récurrence de Levenshtein donne le score contraint.

Pour conserver les mêmes 64 feuilles jusque dans les bandes étroites, une seconde expérience fixe la demi-largeur de chaque paire à \(b_{ij}=\max(b,|m_i-m_j|)\). Les arbres de la matrice exacte et de la matrice approchée utilisent le même ordre de feuilles et la même règle de départage des ex æquo. La distance de Robinson et Foulds~\cite{robinson} compte les groupes internes enracinés pour UPGMA et les bipartitions internes sans racine pour Neighbor-Joining. Pour les ensembles de partitions \(S_1,S_2\), la figure~\ref{fig:sensibilite-arbres} donne
\[
 R=\frac{|S_1\mathbin{\triangle}S_2|}{|S_1|+|S_2|}.
\]
Les longueurs de branches n'interviennent pas dans ce calcul. Les partitions créées par le départage des ex æquo sont conservées.

\chart{biology/sensibilite_arbres.png}{Accessibilité des paires et distance topologique selon la demi-largeur. Le premier panneau compte les paires de la cohorte publique. Le deuxième montre les arbres obtenus en ajustant la bande aux extrémités. Le dernier résume huit jeux synthétiques indépendants de 32 chaînes de même longueur.}{fig:sensibilite-arbres}{scripts/phylogeny/band_sensitivity.py}
À demi-largeur fixe 64, il manque 1\,342 distances sur 2\,016. La bande ajustée rend la matrice complète, avec 74 surestimations. Le RF normalisé vaut \(12/124\) pour UPGMA et \(44/122\) pour Neighbor-Joining. À \(b=128\), les deux topologies retrouvent leur référence alors que huit coûts restent trop élevés en bande ajustée. Sur les témoins synthétiques, le RF normalisé moyen de Neighbor-Joining passe de 0,534 à 0,621 entre les demi-largeurs 0 et 4, même si la bande plus large réduit le coût de chaque paire ou le laisse inchangé. La topologie peut donc s’éloigner de la référence malgré des distances plus précises.

\synthese{Sur cette cohorte, la bande fixe de demi-largeur 323 couvre les 2\,016 paires et restitue les deux topologies. Deux scores dépassent encore leur référence de 5 et 19 éditions. À 512, toutes les distances coïncident.}
